\section{Related Work}
\label{sec:related}

Shortcut learning describes reliance on predictive cues that fail
under a change in the data distribution~\cite{geirhos2020shortcut}.
Eshuijs et al. use planted actor names in movie reviews to study how
attention heads mediate shortcut-driven classification and to
motivate targeted head interventions~\cite{eshuijs2025short}.
This controlled setting motivates our comparison of where a cue
influences predictions and where an intervention can reduce reliance
on it.

The logit lens applies the model's output readout to intermediate
states, revealing how token preferences develop across layers
\cite{nostalgebraist2020logitlens}. A readable preference alone does
not establish which computation causes the final prediction.
Activation patching tests an intervention by replacing selected
activations with those from a reference run. Causal tracing uses
restoration after input corruption to study where factual recall
depends on internal states~\cite{meng2022locating}. These approaches
complement layerwise readouts by measuring changes in model behavior
under controlled interventions.

Deep feature reweighting retrains the final classifier on balanced
data while holding the representation fixed
\cite{kirichenko2023lastlayer}. It provides a precedent for improving
robustness through a restricted parameter update. Activation
Addition instead modifies intermediate states during inference with
a direction estimated from contrasting prompts~\cite{turner2023steering}.
In a video-language setting, Shin et al. trace the weakening of
temporal-order information across layers and reinject an intermediate
temporal signal downstream~\cite{shin2026fades}. Together, these
studies motivate distinguishing where information is readable,
where an intervention affects the output, and which parameters can
be retrained to change a learned behavior.
